\chapter{The Translation Layer}

\epigraph{"I'd rather do math in a general language than do general programming in a mathematical language."}{"New introduction to SciPy" (2013)}

\bigskip

Why do programmers read a mathematician's blog?

Because mathematics provides the abstractions, but code provides the leverage. The blog succeeds because it translates between the two.

The gamma function isn't just a concept; it's a line of Python. The symbol $\lambda$ isn't just an eigenvalue; it's an anonymous function. The Fourier transform isn't just an integral; it's an algorithm. The blog sits at the translation layer, explaining the math to the programmers and the code to the mathematicians.

\bigskip

\section*{Python: The Lingua Franca}

207 posts carry the "Python" tag. Python is the blog's native tongue.

The dominance isn't accidental. Python reads like pseudocode. A snippet explaining the gamma function is accessible to anyone with basic literacy:

\begin{lstlisting}[language=python]
from scipy.special import gamma
print(gamma(5))  # prints 24.0, which is 4!
\end{lstlisting}

The code is almost English. The mathematical content shines through.

The Python posts span the entire history, peaking in 2012 when Python was consolidating its role in science. By the 2020s, Python is assumed. It has become the background radiation of the blog's universe.

\bigskip

\section*{SciPy: The Toolbelt}

80 posts tag "SciPy." This is the library that gives Python its mathematical teeth.

The blog documents the ecosystem not just by teaching it, but by testing it. "Bug in SciPy's erf function" (2010) documents a failure case for complex arguments. This is operational knowledge—the kind gained only by using the tools for real work.

The 2013 post "New introduction to SciPy" makes the thesis explicit: "The Python approach has its advantages—I'd rather do math in a general language than do general programming in a mathematical language."

This is the core argument. Python plus SciPy lets you do mathematics without abandoning general programming. You can read files, process data, and compute integrals in the same script. The cost is complexity; the benefit is power.

\bookfigure{ch05_languages_stacked}{Programming languages over time}

\bookfigure{ch05_languages_bar}{Programming languages by total posts}

\bigskip

\section*{Mathematica: The Symbolic Engine}

74 posts tag "Mathematica."

If Python is for building, Mathematica is for exploring. It handles symbolic manipulation—simplifying expressions, solving equations analytically.

The blog uses Mathematica to check its work. "DFT conventions: NumPy vs Mathematica" (2023) is a classic translation post: it compares how the two systems handle the Discrete Fourier Transform, documenting the subtle differences in convention that trip up practitioners.

Mathematica is the laboratory; Python is the factory. The blog moves between them.

\bigskip

\section*{The Fossils: PowerShell and Perl}

The data preserves extinct interests.

\textbf\{PowerShell\} (47 posts) dominates the early years (2008-2012), shows a lone straggler in 2014, and then disappears.

\textbf\{Perl\} (20 posts) fades similarly, with only sporadic appearances after 2015.

These fossils document the evolution of the field. In 2008, a working consultant needed PowerShell for Windows automation and Perl for text processing. By 2015, Python did both. The blog's topic list is a graveyard of tools that were replaced by better ones.

\bigskip

\section*{Code as Verification}

Many posts include code not just to implement, but to verify.

"Let me check this numerically" is a recurring move. A post will derive a formula, then write a script to test it against a million random cases.

From "10 best rational approximations for pi" (2018):
\begin{quote}
Here are the 10 best rational approximations to $\pi$, found via continued fractions.
\end{quote}

The post includes the code. The code runs. The approximations work.

This is the empiricism of the computer age. Proofs are good; running code is better. The blog insists on both.

\bigskip

\section*{The Bridge}

The blog's value lies in this bridge-building.

A programmer who needs the gamma function can look up the definition (useless) or find a library (opaque). The blog offers the third path: understand the concept \textit\{and\} the implementation.

- Here's the math (gamma extends factorial).
- Here's the code (\texttt\{scipy.special.gamma\}).
- Here's the pitfall (singularities at negative integers).

This is the translation layer. It turns abstract theory into working machinery.
